\documentclass[10pt,a4paper]{amsart}
\usepackage[margin=19mm]{geometry}
\usepackage{amsmath,amssymb,mathtools}
\usepackage[scr=rsfs,cal=euler]{mathalfa}
\usepackage{xfrac}
\usepackage[unicode,hidelinks,bookmarks=false]{hyperref}
\newcommand{\ie}{i.e.\ }
\newcommand{\cf}{cf.\ }
\newcommand{\resp}{respectively\ }
\newcommand{\Subsection}{\subsection}
\providecommand{\nobreakdash}{\nobreak\hbox{-}}
\setlength{\emergencystretch}{2em}
\multlinegap0pt
\usepackage{bm}
\usepackage{bbm}
\usepackage{dsfont}
\usepackage{xspace}
\usepackage{cancel}
\usepackage{mathtools}
\usepackage{cleveref}
\usepackage{amsthm}
\makeatletter
\@ifundefined{definition}{\newtheorem{definition}{Definition}}{}
\@ifundefined{lemma}{\newtheorem{lemma}{Lemma}}{}
\@ifundefined{theorem}{\newtheorem{theorem}{Theorem}}{}
\makeatother
\renewcommand{\Im}[1]{\mathrm{Im}(#1)}
\renewcommand{\Re}[1]{\mathrm{Re}(#1)}
\title[Graph lattice sums and graph zeta functions for long-range i]{Graph lattice sums and graph zeta functions for long-range interacting quantum lattice models}
\author{Andreas Alexander Buchheit, Andreas Rupp}
\date{Source: arXiv:2609.18918v1 (CC BY 4.0)}
\begin{document}
\maketitle

\noindent\emph{Source and licence.} Graph lattice sums and graph zeta functions for long-range interacting quantum lattice models. Version arXiv:2609.18918v1 (CC BY 4.0), distributed under the Creative Commons Attribution 4.0 International licence, as recorded in the arXiv licence metadata for this exact version and shown on the abstract page https://arxiv.org/abs/2609.18918v1. Full rights notice: licenses/arxiv-2609.18918-CC-BY-4.0.txt.\par\smallskip

\noindent\textit{Editorial scope.} Counted unit: the Theorem whose source label is \texttt{thm:multiplication} in the article (source0.tex lines 1142--1164, source label \texttt{thm:multiplication}), delivered together with the article's own complete original proof of it (source0.tex lines 1165--1244, one \texttt{proof} environment of 80 lines). No mathematics is written, repaired or re-derived: statement and proof are the article's own text line for line. Required context retained uncounted: lem:holomorphy (lines 590--620); def:semianalytic (lines 1069--1083). Every definition, display and label of those lines is retained unchanged. Omitted from this package and not redistributed: 1--589, 621--1068, 1084--1141, 1245--1793. Nothing inside a retained excerpt is omitted or altered: every retained body is delivered exactly as the article's lines are written, so no omission receipt of any kind accompanies this package. Citations: the retained excerpts carry no citation command at all, so no bibliography entry of the article is transcribed and no reference is redistributed. Reference adaptation: none was needed and none was made; every reference inside the delivered unit resolves inside it, so no unresolved reference or citation is produced. Printed slips kept literal: no printed word, symbol, hypothesis or number of the counted statement or of its proof was altered. Layout and class: the article's own class is \texttt{amsart} with packages including amsaddr, amssymb, amsthm, thmtools, graphicx, bm, microtype, dsfont, mathtools, tikz, biblatex, csquotes, comment, hyperref; the wrapper is the lane's pinned \texttt{amsart} editorial wrapper at 10pt on A4, which restates the theorem-like environments the retained text needs and adds the article's own macro definitions that the retained text needs (\texttt{\textbackslash Im}, \texttt{\textbackslash Re}), with extra wrapper packages (bm, bbm, dsfont, xspace, cancel, mathtools, cleveref). The delivered numbering is the unit's own. Assets: no retained span contains a figure environment, a graphics-inclusion command, a PDF-page inclusion, a table or a TikZ picture; no image file is copied into this package, the package declares no asset dependency and no image file is redistributed with it. Licence: version arXiv:2609.18918v1 of this article is distributed under the Creative Commons Attribution 4.0 International licence, as recorded in the arXiv licence metadata for this exact version and shown on the abstract page https://arxiv.org/abs/2609.18918v1. A full rights notice is delivered at licenses/arxiv-2609.18918-CC-BY-4.0.txt. Characters: every retained excerpt is pure ASCII, so the delivered engine, XeLaTeX with its default Latin Modern text font, reports no missing character.
\begin{lemma}[Holomorphy and singularity structure of the Epstein zeta function]
    \label{lem:holomorphy}
    Define the set $D_L\subseteq \mathds C^d$ as the  following intersection of $d$-dimensional complex cones with origins at $L\subseteq \mathds R^d$,
$$
D_L=\{\bm u\in\mathds C^d:|\Re{\bm u}-\bm z|>|\Im{\bm u}|\ \forall \bm z\in L\}.
$$
\begin{enumerate} \item The Epstein zeta function can be holomorphically extended to  
$$
(\nu,\bm k)\in 
\mathds C\times D_{\Lambda^*}.
$$
For $\bm k\in\Lambda^*$, the Epstein zeta function is holomorphic in $\nu\in\mathds C\setminus\{d\}$ with a simple pole in $\nu=d$.  

\item The singularity of the Epstein zeta function at $\bm k=\bm 0$ can be isolated as
    \[
    Z_{\Lambda,\nu}(\bm k) = Z_{\Lambda,\nu}^{\mathrm{reg}}(\bm k)+\hat s_\nu(\bm k)/V_\Lambda,
    \]
    where $Z_{\Lambda,\nu}^{\mathrm{reg}}$ is analytic on $(\mathds R^d\setminus \Lambda^\ast)\cup \{\bm 0\}$ and where $\hat s_\nu(\bm k)$ is the distributional Fourier transform of $s_\nu(\bm x)=\vert \bm x\vert^{-\nu}$, which reads for $\nu\in \mathds C\setminus (2\mathds N +d)$,
    \[
    \hat s_\nu(\bm k) = c_\nu \vert \bm k\vert^{\nu-d},\qquad c_\nu=\pi^{\nu-d/2}\frac{ \Gamma ((d-\nu)/2)}{\Gamma(\nu/2) }.
    \]
    In case that $\nu=d+2n$, $n\in \mathds N$, the Fourier transform is uniquely defined up to a polynomial of order $2n$. We adopt the choice,
\[\hat s_{d+2n}(\bm y)= \frac{\pi^{n+d/2}}{\Gamma(n+d/2)}\frac{(-1)^{n+1}}{n!} ( \pi \bm k^2 )^{n} \log (\pi  \bm k^{2}).
    \]
\item For $\nu\in d+2 \mathds N_0$, the regularised Epstein zeta function can be holomorphically extended to  $\bm k \in D_{\Lambda^\ast\setminus\{\bm 0\}}$,
and, outside these particular $\nu$ values, the funtion extends holomorphically to
\[
(\nu,\bm k) \in \big(\mathds C\setminus (d+2\mathds N_0)\big)\times D_{\Lambda^\ast\setminus\{\bm 0\}}.
\]
\end{enumerate}
\end{lemma}  

\begin{definition}[Semi-analytic kernels and graph lattice sums]
\label{def:semianalytic}
We call a kernel $K\in\ell^1(\Lambda)$ semi-analytic if it admits a
representation
\[
K(\bm x)=a(\bm x)+\sum_{j=1}^{N}b_j\,\mathcal K_{\nu_j}(\bm x),
\qquad N\in\mathds N_0,
\]
with $b_j,\nu_j\in\mathds C$, $\mathrm{Re}(\nu_j)>d$, and even coefficients $a(\bm x)\in\mathds C$ obeying the bound $|a(\bm x)|\le C|\bm x|^{-(d+2)}$ for all $\bm x\in\Lambda\setminus\{\bm 0\}$ and some $C>0$. We note that this representation is not unique and operations within this section will act on a given representation. Finally, a graph lattice sum $\mathcal Z_G$ is called semi-analytic, if it is the lattice Fourier transform of a semi-analytic kernel, namely
\[
\mathcal Z_G(\bm k)=\sum_{\bm x\in\Lambda}a(\bm x)e^{-2\pi i\bm x\cdot\bm k}
+\sum_{j=1}^{N}b_j\,Z_{\Lambda,\nu_j}(\bm k),
\]
with $Z_{\Lambda,\nu}$ the Epstein zeta function.
\end{definition}

\begin{theorem}[Multiplication of semi-analytic graph lattice sums]
\label{thm:multiplication}
Let $\mathcal Z_{G_1},\mathcal Z_{G_2}$ be semi-analytic as in \cref{def:semianalytic}, with representations $(a_\eta,\bm b_\eta,\bm\nu_\eta)$ and $\bm b_\eta,\bm\nu_\eta\in\mathds C^{N_\eta}$, $\eta\in\{1,2\}$, where $\mathrm{Re}(\nu_{i\eta})>d$ and $\nu_{i1},\nu_{j2}\notin d+2\mathds N_0$. Let the regular parts satisfy $|a_\eta(\bm x)|\le C_\eta|\bm x|^{-\gamma_\eta}$ on $\Lambda\setminus\{\bm 0\}$ with $\gamma_\eta>d$. Then $\mathcal Z_{G_1}\mathcal Z_{G_2}$ is semi-analytic, and
\[
\begin{aligned}
\mathcal Z_{G_1}(\bm k)\,\mathcal Z_{G_2}(\bm k)
&=\sum_{\bm x\in\Lambda}a_3(\bm x)\,e^{-2\pi i\bm x\cdot\bm k}+\sum_{i=1}^{N_1}\sum_{j=1}^{N_2} b_{i1}b_{j2}
\frac{c_{\nu_{i1}}\,c_{\nu_{j2}}}{V_\Lambda\,c_{\nu_{i1}+\nu_{j2}-d}}\,
Z_{\Lambda,\nu_{i1}+\nu_{j2}-d}(\bm k)\\
&\quad+\mathcal Z_{G_2}(\bm 0)\sum_{i=1}^{N_1}b_{i1}\,Z_{\Lambda,\nu_{i1}}(\bm k)
+\mathcal Z_{G_1}(\bm 0)\sum_{j=1}^{N_2}b_{j2}\,Z_{\Lambda,\nu_{j2}}(\bm k),
\end{aligned}
\]
where $1/c_{\nu_{i1}+\nu_{j2}-d}$ is set to $0$ when
$\nu_{i1}+\nu_{j2}-d\in d+2\mathds N_0$. The exponents of the Epstein terms again satisfy
$\mathrm{Re}(\nu_{i1}+\nu_{j2}-d)>d$. The regular part $a_3$ is even and obeys the bound
$|a_3(\bm x)|\le C_3|\bm x|^{-\gamma_3}$ on $\Lambda\setminus\{\bm 0\}$ for some
$C_3>0$, with
\[
\gamma_3=\min\Big(\gamma_1,\;\gamma_2,\;\min_{i}\mathrm{Re}(\nu_{i1})+2,\;\min_{j}\mathrm{Re}(\nu_{j2})+2\Big),
\]
with minima over empty index sets taken as $\infty$. In particular, $\gamma_3\ge d+2$ if $\gamma_1,\gamma_2\ge d+2$, so that semi-analyticity from \cref{def:semianalytic} is preserved.
\end{theorem}

\begin{proof}
We begin with a preparatory step, decomposing each Epstein zeta function in the two graph lattice sums into its regular part and a Riesz kernel,
$Z_{\Lambda,\nu}=Z^{\mathrm{reg}}_{\Lambda,\nu}+\hat s_\nu/V_\Lambda$ with
$\hat s_\nu(\bm k)=c_\nu|\bm k|^{\nu-d}$ the distributional Fourier transform of $\vert\bm \cdot\vert^{-\nu}$ on $\mathds R^d$, see \cref{lem:holomorphy}. Note that $Z^{\mathrm{reg}}_{\Lambda,\nu}$ is analytic in a complex
neighbourhood of $\bm k=\bm 0$ and $Z_{\Lambda,\nu}$ is analytic on
$\mathrm{BZ}\setminus\Lambda^\ast$. We can now separate each graph lattice sum $\mathcal Z_{G_\eta}$ into singular terms at $\bm k=\bm 0$ plus a remainder, 
\[
\mathcal Z_{G_\eta}=\frac{1}{V_\Lambda}S_\eta+R_\eta
\]
with
\[
S_\eta=\sum_{i=1}^{N_\eta} b_{i\eta}\,\hat s_{\nu_{i\eta}},
\qquad
A_\eta=\sum_{i=1}^{N_\eta} b_{i\eta}\,Z^{\mathrm{reg}}_{\Lambda,\nu_{i\eta}},
\qquad
R_\eta=A_\eta+\hat a_\eta,
\]
with the Fourier series $\hat a_\eta(\bm k)= (\mathcal F_\Lambda a_\eta)(\bm k)$.
Note that both $A_\eta$ and $\hat a_\eta$ are even, the former as $-\Lambda=\Lambda$, the latter by semi-analyticity. Therefore, also $R_\eta$ is even.

In the next step, fix a radially symmetric cutoff function $\chi\in C^\infty(\mathrm{BZ},\mathds R_{\ge0})$ with $\chi=1$ on $|\bm k|<r/2$ and $\chi=0$ on $|\bm k|>r$. Choose the radius $r>0$ small enough, such that both $r<\rho_{\Lambda^\ast}$ with $\rho_{\Lambda^\ast}=\min_{\bm q\in\Lambda^\ast\setminus\{\bm 0\}}|\bm q|$, and such that $\chi$ vanishes in an environment of $\partial\mathrm{BZ}$.  By \cref{lem:holomorphy}, $Z^{\mathrm{reg}}_{\Lambda,\nu_{i\eta}}$ is then analytic on a complex ball of radius
$(\rho_{\Lambda^\ast}-r)/\sqrt2$ about every point of $\mathrm{supp}\,\chi$. Split the product of the two graph lattice sums into $\mathcal Z_{G_1} \mathcal Z_{G_2}=\mathcal Z_{G_1} \mathcal Z_{G_2}(1-\chi)+\mathcal Z_{G_1} \mathcal Z_{G_2}\chi$, the first term periodic and as regular as $\hat a_\eta$, the second carrying the localized power-law singularities at $\bm k=\bm 0$.

(\emph{Step 1:\,Regular part}) We begin the discussion with the regular part. Expand $\mathcal Z_{G_1}\mathcal Z_{G_2}(1-\chi)$ into the four products of Epstein and Fourier parts. The term involving products of Epstein zeta functions is a product of analytic functions on $\mathrm{supp}(1-\chi)$ multiplied by a smooth periodic cutoff. 
Therefore, it is smooth and periodic on all of $\mathrm{BZ}$, thus exhibiting superalgebraically decaying Fourier coefficients. 

The Fourier coefficients of the two mixed terms are given by the lattice convolution of a superalgebraically decaying function with $a_\eta$, hence decaying asymptotically as $|\bm x|^{-\gamma_\eta}$. 
Finally, the Fourier coefficients of $\hat a_1 \hat a_2$ are $\sum_{\bm y\in \Lambda}a_1(\bm y)a_2(\bm x-\bm y)$. Split the sum at $|\bm y|=|\bm x|/2$. For $|\bm y|\le |\bm x|/2$, bound $|a_2(\bm x-\bm y)|\le C_2(|\bm x|/2)^{-\gamma_2}$ by the inverse triangle inequality. For $|\bm y|>|\bm x|/2$, bound $|a_1(\bm y)|\le C_1 (|\bm x|/2)^{-\gamma_1}$. The absolute value of the sum is bounded by the $\ell^1$ norm of the remaining summand. This yields
\[
\sum_{\bm y\in \Lambda}|a_1(\bm y) a_2(\bm x-\bm y)|\le \Vert a_1\Vert_{\ell^1}C_2(|\bm x|/2)^{-\gamma_2}+\Vert a_2\Vert_{\ell^1}C_1(|\bm x|/2)^{-\gamma_1}\le C |\bm x|^{-\min(\gamma_1,\gamma_2)},
\]
for some $C>0$. All of this is absorbed into $a_3$.

(\emph{Step 2:\,Singular part}) We then proceed with the localized singular part. On $\mathrm{supp}\,\chi$,
\[
\mathcal Z_{G_1}\mathcal Z_{G_2}
=\frac{S_1S_2}{V_\Lambda^{2}}
+\frac{S_1R_2+S_2R_1}{V_\Lambda}
+R_1R_2,
\]
yielding three groups requiring separate analysis. 
For the first group, we have
\[
\frac{S_1S_2}{V_\Lambda^{2}}=\sum_{i=1}^{N_1}\sum_{j=1}^{N_2}b_{i1}b_{j2} \frac{\hat s_{\nu_{i1}}\hat s_{\nu_{j2}}}{V_\Lambda^{2}}.
\]
For $\nu\notin d+2\mathds N_0$, we have $\hat s_\nu=c_\nu|\bm k|^{\nu-d}$, avoiding potential logarithmic terms. By the restriction on $\nu_{i\eta}$, the resulting product is therefore again a power-law,
\[
\frac{\hat s_{\nu_{i1}}\hat s_{\nu_{j2}}}{V_\Lambda^{2}}=\frac{c_{\nu_{i1}}c_{\nu_{j2}}}{V_\Lambda^{2}}\,|\bm k|^{\nu_{i1}+\nu_{j2}-2d}=\frac{c_{\nu_{i1}}c_{\nu_{j2}}}{V_\Lambda\,c_{\nu_{i1}+\nu_{j2}-d}}\frac{\hat s_{\nu_{i1}+\nu_{j2}-d}}{V_\Lambda}.
\]
After writing $\hat s_\mu/V_\Lambda=Z_{\Lambda,\mu}-Z^{\mathrm{reg}}_{\Lambda,\mu}$, we note that $Z^{\mathrm{reg}}_{\Lambda,\mu}\chi$ is smooth and periodic, which can therefore be absorbed in the Fourier series. This provides the double sum of $N_1N_2$ Epstein zeta functions. Here and below, $Z_{\Lambda,\mu}\chi$ differs from $Z_{\Lambda,\mu}$ by the smooth periodic function $Z_{\Lambda,\mu}(1-\chi)$, which is absorbed into the Fourier series as well. If $\nu_{i1}+\nu_{j2}-d=d+2n$, the product is a polynomial in $\bm k$ and hence analytic, and its product with $\chi$ can therefore be completely absorbed into the Fourier series. Correspondingly, the coefficient $c_{\nu_{i1}+\nu_{j2}-d}$ exhibits a pole there, so that with the convention $1/c_{\nu_{i1}+\nu_{j2}-d}=0$ no Epstein term is generated.

For the second group, $\mathrm{Re}(\nu)>d$ gives $\hat s_\nu(\bm 0)=0$, yielding
$R_\mu(\bm 0)=\mathcal Z_{G_\mu}(\bm 0)$. For $\eta\neq\mu$, apply the split
\[
S_\eta R_\mu=\mathcal Z_{G_\mu}(\bm 0)\,S_\eta
+S_\eta\big(A_\mu-A_\mu(\bm 0)\big)
+S_\eta\big(\hat a_\mu-\hat a_\mu(\bm 0)\big).
\]
The first term contributes
$\mathcal Z_{G_\mu}(\bm 0)\sum_i b_{i\eta}Z_{\Lambda,\nu_{i\eta}}$. For $(\eta,\mu)=(1,2)$ and $(2,1)$, we obtain the two single sums in the theorem statement. In the second term, we use that $A_\mu$ is analytic and even. Hence, its
first-order Taylor coefficient vanishes and Cauchy estimates on the ball of radius
$(\rho_{\Lambda^\ast}-r)/\sqrt2$ give $A_\mu-A_\mu(\bm 0)=\mathcal O(|\bm k|^{2})$ on $\mathrm{supp}\,\chi$. Each summand is then $|\bm k|^{\nu_{i\eta}-d}\chi$ times a smooth function vanishing to second order at $\bm 0$. Expanding this function into homogeneous polynomials shows that the Fourier coefficients of the summand are of order $|\bm x|^{-(\mathrm{Re}(\nu_{i\eta})+2)}$.  The third term needs to be estimated on the coefficient side. Up to superalgebraically decaying corrections, the coefficients of $\hat s_\nu\chi/V_\Lambda$ are equal to $\mathcal K_\nu$, and $\hat a_\mu(\bm 0)=\sum_{\bm y\in \Lambda}a_\mu(\bm y)$. Therefore, the coefficients of $S_\eta(\hat a_\mu-\hat a_\mu(\bm 0))\chi/V_\Lambda$ are, up to corrections of order $|\bm x|^{-\gamma_\mu}$,
\[
\sum_{i=1}^{N_\eta} b_{i\eta}\sum_{\bm y\in\Lambda}
\Big(\mathcal K_{\nu_{i\eta}}(\bm x-\bm y)-\mathcal K_{\nu_{i\eta}}(\bm x)\Big)
a_\mu(\bm y).
\]
As before, split the inner sum at $|\bm y|=|\bm x|/2$. On $|\bm y|\le|\bm x|/2$, the terms in brackets obey
\[\mathcal K_{\nu_{i\eta}}(\bm x-\bm y)-\mathcal K_{\nu_{i\eta}}(\bm x)=-\bm y\cdot\nabla\mathcal K_{\nu_{i\eta}}(\bm x)
+\mathcal O\big(|\bm y|^{2}|\bm x|^{-\mathrm{Re}(\nu_{i\eta})-2}\big).\] The linear term cancels because the summation region is symmetric and $a_\mu$ is even, and $\sum_{|\bm y|\le R}|\bm y|^{2}|a_\mu(\bm y)| =\mathcal O\big(C_\mu R^{(d+2-\gamma_\mu)_+}\big)$, with a logarithm in place of the power in case that $\gamma_\mu=d+2$. 
This region therefore contributes the exponent $\mathrm{Re}(\nu_{i\eta})+\min(2,\gamma_\mu-d)$. On $|\bm y|>|\bm x|/2$ the two terms are bounded separately, by $C_\mu(|\bm x|/2)^{-\gamma_\mu}\|\mathcal K_{\nu_{i\eta}}\|_{\ell^1}$ and by $|\bm x|^{-\mathrm{Re}(\nu_{i\eta})}\sum_{|\bm y|>|\bm x|/2}|a_\mu(\bm y)|$, of exponents $\gamma_\mu$ and $\mathrm{Re}(\nu_{i\eta})+\gamma_\mu-d$ respectively, where $\|\mathcal K_\nu\|_{\ell^1}=Z_{\Lambda,\mathrm{Re}(\nu)}(\bm 0)$.

For the third group, $A_1A_2\chi$ is smooth and periodic, $A_\eta \hat a_\mu\chi$ has
coefficients given by a convolution of a superalgebraically decaying sequence with
$a_\mu$, and $\hat a_1 \hat a_2$ was treated in Step~1. Therefore, all decay as
$|\bm x|^{-\min(\gamma_1,\gamma_2)}$ or faster.

Taking the minimum over all exponents yields $\gamma_3$. Step~1 and the third group of Step~2 give $\min(\gamma_1,\gamma_2)$. In the second group of Step~2, the second term gives $\min_i\mathrm{Re}(\nu_{i\eta})+2$, and the third term gives $\gamma_\mu$ and $\mathrm{Re}(\nu_{i\eta})+\min(2,\gamma_\mu-d)$. The latter equals $\mathrm{Re}(\nu_{i\eta})+2$ for $\gamma_\mu\ge d+2$ and exceeds $\gamma_\mu$ otherwise. All other terms are smooth and periodic. The logarithm at $\gamma_\mu=d+2$ does not enter the bound, since $|\bm x|^{-\mathrm{Re}(\nu_{i\eta})-2}\log|\bm x|$ decays faster than $|\bm x|^{-\gamma_\mu}$.
Finally, every contribution is a product or convolution of even functions, as $\chi$ is radial, so $a_3$ is even and $\mathcal Z_{G_1}\mathcal Z_{G_2}$ is semi-analytic in the sense of \cref{def:semianalytic}.
\end{proof}

\end{document}
